\subsection*{Croqtile 值编号简介}

在 Croqtile 中，值编号过程用来支撑形状推导：要么直接推导出编译时已知的形状，要么生成形状计算表达式。

许多值编号方法着眼于消除冗余计算，也就是那些始终产生相同值的表达式；Croqtile 的值编号则想建立一种渐进的方法，化简自己的领域专用构造 \co@mdspan@。在这种语境下，常量传播、代数化简等技术更为关键。

实现上，Croqtile 采用一种简单的\textbf{基于作用域}的值编号方法：它是局部值编号的扩展，配上一张词法作用域的值编号表[1]。与其他方法不同，Croqtile 不需要构造基本块、剪枝 SSA，也不需要编译器后端，因为它直接把值编号施加在抽象语法树上。但为了配合这种做法，必须给 Croqtile 的语法规则加上若干限制，尤其是防止值被改写。
